\begin{abstract}
\textbf{Background:} Removing covariates from clinical EEG features can create spurious discrimination in nonlinear classifiers. For windows nested within patients, the size of this effect and the error frequencies of a feature-shuffle audit were not reported in the references reviewed.
\textbf{Objective:} To measure the discrimination created by covariate adjustment when the EEG carries no label information, and the error frequencies of rules detecting it, in patient-level epilepsy classification.
\textbf{Methods:} We analyzed first sessions of 200 patients (100 with epilepsy) in the TUH EEG Epilepsy Corpus v2.0.1 (not a first-seizure population) on 50 patient-level folds. A three-view fusion pipeline (F2) was audited with null experiments (EEG permuted between patients), simulations, a plasmode (real EEG features, covariate-generated labels), null and positive rule controls, non-EEG baselines, matched strata, and a selection-optimism correction.
\textbf{Results:} A reference fitted on training controls leaked (null AUC 56.4); in the plasmode, adjustment created up to 11.9 percentage points (pp) of spurious discrimination in nonlinear classifiers and at most 1.0 pp in logistic regression (independent cohort: 12.7 and 0.3 pp). With 20 replicates, the absolute rule declared leakage in 0.1\% of leakage-free draws and 88\% for the leaking reference, but the paired rule declared leakage in 32--48\% of null draws for raw tree pipelines (2.7--4.2\% with 50 replicates). In simulation, the leak was not shown to grow with within-patient window correlation ($-$1.8 pp); in one simulated cell, aggregating to patient means first avoided it only at low correlation. F2 (per-fold AUC 69.7; 95\% CI 61.0--78.5) was not shown to exceed the acquisition-metadata baseline (65.5) or a model without adjustment (67.5); EEG added 7.6 pp to all non-EEG predictors ($p = 0.082$).
\textbf{Conclusion:} The audit can flag adjustment-created discrimination, but rank-normal mitigation depends on the classifier and rules need enough replicates. Evidence comes from simulations and two archives, one over 50 rounds.
\end{abstract}
\noindent\textbf{Keywords:} electroencephalography; epilepsy; data leakage; confounder; covariate adjustment; predictive model validation
